\section{문제 정의}
\label{sec:problem}

\begin{revision}

장면의 사건을 $E_s=(e_0,\ldots,e_{n-1})$로 두고 자료상 연속성이 확인된 구간에서만 의무를 이어받는다. 사건은 행동, 새 의무, 조건별 증거와 원문 구절을 가진다. 의무 $o$는 ID, 대상과 해제 조건 ID $c(o)$에 연결된다. 조건 증거 $q_i(c)\in\{\mathsf{T},\mathsf{F},?\}$에서 참은 해당 조건 충족, 거짓은 미충족 확인, $?$는 알 수 없음을 뜻한다. 새 행동 근거 B는 ID가 다른 A의 해제 증거가 아니다.
예를 들어 의무 A의 시작 사건은 정지 행동과 ``보행자가 통과할 때까지''라는 해제 조건을 함께 갖는다. 이후 사건의 녹색 신호 B는 진행 근거로 기록할 수 있지만, 보행자 조건 A의 증거를 채우려면 같은 조건에 대응하는 구절이 있어야 한다. 사건 간 거리가 길어져도 A와 B를 바꿔 연결하지 않으며, 중간 사건에 A의 증거가 없다는 이유만으로 A를 지우지 않는다. 이처럼 검사 단위는 행동 단어 한 개가 아니라 의무의 시작--조건 유지--해제 증거--후속 행동의 연결이다.


\subsection{주석 문제와 검사기 처리 오류의 분리}
\label{sec:error-types}
P1--P3은 입력의 행동ㆍ요구 관계에 관한 문제이고 P4는 검사기의 상태 처리 문제다. P5는 별도의 속도 반응 시간 검사다. 기존 40쌍을 이 구분으로 재집계한다.

\begin{align}
C_1(i)&=\mathrm{GO}_i\land\exists o\in A_i^-:q_i^*(c(o))=\mathsf{F},\label{eq:hold-conflict}\\
C_2(i)&=\exists o:\mathrm{release}_i(o)\land q_i(c(o))=\mathsf{F},\\
C_3(i)&=\mathrm{GO}_i\land p_i^*=\mathsf{F},\\
F_4(i)&=\exists o\in A_{i-1}:q_i^*(c(o))=\mathsf{T}\land o\in A_i,\\
M(i)&=\mathrm{GO}_i\land\exists o\in A_i^-:q_i^*(c(o))=? .\label{eq:missing-evidence}
\end{align}

$A_i^-$는 새 의무와 해제 증거를 반영한 뒤 현재 행동 직전의 활성 의무 집합이다. $q_i^*$는 유효한 최근 증거를, $p_i^*$는 저장된 선행조건 증거를 뜻한다. P1은 미해제가 확인된 의무와 진행의 충돌, P2는 미충족 조건을 충족했다고 처리한 조기 해제, P3은 선행조건 미충족 상태의 진행이다. $F_4$는 해제가 확인되었는데도 해당 의무를 상태에 남기는 처리 오류다. $M$은 진행 시 해제 근거 부족이며 확정 모순에 포함하지 않는다.

기존 검사기의 P2 출력은 진행 사건의 \texttt{release=false}에서, P4 출력은 같은 대상의 정지ㆍ양보 반복과 \texttt{release=true}가 함께 주어진 시험에서 발생한다. 이 구현상 표시는 실제 사람이 조기 해제했다고 독립 확인한 결과가 아니다. P4 변형 10건은 의도적으로 만든 처리 시험으로 보고한다. 보완 모델은 별도의 해제 명령이나 P4 고장 주입을 받지 않고 조건 ID의 참 증거로 정상 상태를 갱신한다. 따라서 보완 모델의 144/144를 P1--P4 전체의 새 성능이라고 서술하지 않는다.

\subsection{조건별 갱신과 판정 의미}
\label{sec:semantics}
보완 모델의 갱신은 다음과 같다. 미상 입력은 최근의 알려진 증거를 덮지 않는다.
\begin{equation}
q_i^*(c)=\begin{cases}q_i(c),&q_i(c)\ne ?,\\q_{i-1}^*(c),&q_i(c)=?,\end{cases}
\end{equation}
\begin{equation}
A_i^-=(A_{i-1}\cup O_i)\setminus\{o:q_i^*(c(o))=\mathsf{T}\},\qquad A_i=A_i^- .
\label{eq:active-update}
\end{equation}
이 보존 규칙은 같은 대상ㆍ조건이 유지되고 증거의 유효 구간이 계속된다는 작성 자료의 가정에 의존한다. 실제 객체의 추적ㆍ증거 만료ㆍ재등장은 모델링하지 않는다. 한 번 해제된 의무는 같은 ID의 재생성 입력이 없으면 다시 활성화하지 않는다.

정지나 대기 중 해제 여부가 미상이라는 이유만으로 문제를 기록하지 않는다. 진행할 때 남은 의무가 거짓이면 모순, 미상이면 근거 부족이다. 복수 의무는 각각의 조건을 확인하며 하나의 해제가 다른 의무를 지우지 않는다. 같은 사건에 모순과 근거 부족이 함께 있으면 모순을 최종 판정으로 우선하되 각각의 기록을 보존한다. 전체 판정은 누적 모순이 있으면 모순, 그렇지 않고 진행 중 근거 부족이 있었으면 미상, 둘 다 없으면 표현한 의무와 양립이다. 마지막 판정은 모든 위험이 없다는 뜻이 아니다.

장면 끝에서 미래의 해제를 보지 못했다는 사실만으로 모순을 추가하지 않는다. 미상 상태에서 진행하지 않은 대기열은 현재 검사 범위에서 문제 없음으로 끝날 수 있다. 근거 없이 다음 장면과 연결하지 않으며 다음 장면은 초기화한다. 임의의 미래 완료 의무나 제한 시간을 추가한 liveness 명세는 이번 중심 검사에 없다.\subsection{최초 문제 위치와 별도 반응 질문}
\label{sec:issue-definition}
조건별 중심 모델에서 확정 모순의 최초 사건과 근거 부족의 최초 사건은 따로 기록한다.
\begin{equation}
i_{\mathrm{bad}}=\min\{i:C_1(i)\lor C_3(i)\},\quad
i_{\mathrm{miss}}=\min\{i:M(i)\}.
\end{equation}
해당 집합이 비면 그 최초 위치는 없다. 최종 모순 판정에는 $i_{\mathrm{bad}}$를, 모순 없이 근거 부족만 있으면 $i_{\mathrm{miss}}$를 대표 문제 위치로 사용한다. 뒤의 정상 해제가 앞서 근거 없이 진행한 사건을 없애지는 않는다. 이는 현재 순간의 활성 의무와 전체 경로의 문제 기록을 구분하는 이유다. P2/P4 출력은 앞서 설명한 기존 회귀 경로에서만 별도 다룬다.

보조 질문 P5는 사건 시각 $\tau_i$, 요구 행동 $a_i$, 선택한 제한 시간 $D$와 속도 반응 조건 $R_{a_i}$에 대해
\begin{equation}
\exists t_r\in[\tau_i,\tau_i+D):R_{a_i}(t_r)
\label{eq:response-contract}
\end{equation}
를 확인한다. 기록 속도에서 찾은 $t_r$은 A의 객체 해제를 직접 확인하는 증거가 아니다. 따라서 P5는 조건별 중심 모델의 정량 시계 속성이 아니라 기록 반응에 대한 별도 분석이다. 문장열이 명시된 의무와 양립하는지와 자차 반응이 선택한 시간 설정에 맞는지는 서로 다른 질문이다.


\end{revision}
